\subsection{ভগ্নাংশ, মিশ্র ভিত্তি ও ধারসহ উদাহরণ}
ভগ্নাংশ থাকলে সংখ্যাগুলো \textbf{র‍্যাডিক্স পয়েন্ট বরাবর} সাজাতে হয়। প্রয়োজনীয় শূন্য পূর্ণাংশের বামে বা ভগ্নাংশের শেষে বসানো যায়।

\chThreeExample{১}{ভগ্নাংশের যোগ ও বিয়োগ}
\chThreeAnswer
\[\begin{array}{r}1011.101\\+0110.011\\\hline10010.000\end{array}
\qquad
\begin{array}{r}1011.101\\-0110.011\\\hline0101.010\end{array}\]
যাচাই: $11.625+6.375=18$ এবং $11.625-6.375=5.25$। ভগ্নাংশের যোগে পাওয়া ক্যারি পূর্ণাংশে যায়; বিয়োগেও বিন্দুর দুই পাশে ধার চলতে পারে।

\chThreeExample{২}{একাধিক শূন্য অতিক্রম করে ধার}
$(10000)_2-(1)_2$।
\chThreeAnswer পাঁচ বিটে $10000-00001=01111$। বাম প্রান্তের ১ থেকে ধার আসার সময় মাঝের শূন্যগুলো ধার দিয়ে ১ হয়ে থাকে; ডানের ঘরে $10_2-1=1$ হয়। দশমিকে $16-1=15$।

\chThreeExample{৩}{বই ও কলমের মূল্য}
একটি বইয়ের মূল্য $(A6)_{16}$ টাকা, কলমের মূল্য $(35)_8$ টাকা। মোট মূল্য অক্টালে লেখো।
\chThreeAnswer $A6=166$, $(35)_8=29$। মোট $195$ টাকা।
\[195\div8=24\;R3,\quad24\div8=3\;R0,\quad3\div8=0\;R3.\]
অতএব মোট $(303)_8$ টাকা। সরাসরি যোগ করতে চাইলে প্রথমে $(A6)_{16}=(246)_8$, তারপর $(246)_8+(35)_8=(303)_8$।

\chThreeExample{৪}{অক্টাল ও হেক্সাডেসিমাল ভগ্নাংশ}
\chThreeAnswer
\[(17.6)_8+(2.3)_8=(22.1)_8,\qquad (2A.8)_{16}-(B.C)_{16}=(1E.C)_{16}.\]
প্রথমটিতে $6+3=9_{10}=11_8$, তাই ভগ্নাংশে ১ রেখে পূর্ণাংশে ক্যারি ১। দ্বিতীয়টিতে $8-C$ করতে পূর্ণাংশ থেকে ধার নিয়ে $8+16-12=12=C$ হয়; পূর্ণাংশে অবশিষ্ট $29_{16}-B_{16}=1E_{16}$।

\chThreeExample{৫}{ভিন্ন ভিত্তির ভগ্নাংশ বাইনারিতে যোগ}
$(2B.4)_{16}$ ও $(15.6)_8$ বাইনারিতে যোগ করে ফল হেক্সাডেসিমালে প্রকাশ করো।
\chThreeAnswer
\[(2B.4)_{16}=(101011.01)_2,\quad (15.6)_8=(1101.11)_2.\]
\[\begin{array}{r}101011.01\\+001101.11\\\hline111001.00\end{array}\]
অতএব ফল $(39)_{16}$; যাচাই $43.25+13.75=57=3\times16+9$।

\smallhead{সম্পূরক অনুশীলন: বাইনারি গুণ ও ভাগ}
যোগ-বিয়োগ বোঝার পরে শিফটের প্রয়োগ হিসেবে নিচের ছোট উদাহরণগুলো দেখা যায়। গুণে গুণকের বিট ১ হলে গুণ্য লেখো, ০ হলে শূন্য; পরের বিটের আংশিক গুণফল এক ঘর বামে সরাও।
\[(1011)_2(101)_2=(1011)_2+(101100)_2=(110111)_2.\]
দশমিকে $11\times5=55$। ভগ্নাংশে আগে বিন্দু সরিয়ে পূর্ণসংখ্যার মতো গুণ করো; শেষে দুই সংখ্যার মোট ভগ্নাংশ বিট অনুযায়ী বিন্দু বসাও।
\[(1.1)_2(10.1)_2=(11)_2(101)_2\div2^2=(1111)_2\div4=(11.11)_2.\]
যাচাই $1.5\times2.5=3.75$। ভাগে দীর্ঘ ভাগের নিয়ম ব্যবহার হয়: $(11001)_2\div(101)_2=(101)_2$, কারণ $25\div5=5$। ভাজক শূন্য হতে পারে না।

% Mahbub teaching-type follow-up, 2026-10-14
\chThreeExample{৬}{ভিন্ন ভিত্তির ভগ্নাংশ বিয়োগ}
$(3C.8)_{16}$ থেকে $(26.75)_{10}$ বিয়োগ করে ফল বাইনারিতে লেখো।
\chThreeAnswer $(3C.8)_{16}=60.5=(111100.10)_2$ এবং $26.75=(11010.11)_2$।
\[\begin{array}{r}111100.10\\-011010.11\\\hline100001.11\end{array}\]
ফল $33.75=(100001.11)_2=(21.C)_{16}$। বিন্দুর পর ০ থেকে ১ বিয়োগ করতে ধার পূর্ণাংশ পর্যন্ত যায়। দশমিক যাচাই $60.5-26.75=33.75$।

\chThreeExample{৭}{তিনটি ভিন্ন ভিত্তির পরিমাণের মোট ও তুলনা}
তিন শিক্ষার্থী একটি প্রকল্পে যথাক্রমে $(101101)_2$, $(36)_8$ ও $(2A)_{16}$ টাকা দিয়েছে। মোট কত এবং কে সবচেয়ে বেশি দিয়েছে?
\chThreeAnswer পরিমাণগুলো দশমিকে ৪৫, ৩০ ও ৪২। বাইনারিতে যোগ:
\[\begin{array}{r}0101101\\0011110\\+0101010\\\hline1110101\end{array}\]
মোট $117=(1110101)_2=(165)_8=(75)_{16}$ টাকা। প্রথম শিক্ষার্থী সর্বাধিক দিয়েছে; তৃতীয়জনের চেয়ে $45-42=3$ টাকা বেশি। তুলনায় অঙ্কের দৈর্ঘ্য নয়, একই ভিত্তিতে প্রকাশিত মান বিবেচ্য।

\chThreeExample{৮}{মিশ্র ভিত্তির যোগে অজানা সংখ্যা}
$(3C)_{16}+(x)_8=(91)_{16}$ হলে $(x)_8$ নির্ণয় করো।
\chThreeAnswer $91_{16}=145$ এবং $3C_{16}=60$। তাই $x$-এর মান $145-60=85$। $85\div8=10\;R5$, $10\div8=1\;R2$, $1\div8=0\;R1$; অতএব $(x)_8=(125)_8$। যাচাই: $60+(64+16+5)=145$। এখানে অজানা সংখ্যার অঙ্কগুলো বের করা হয়েছে; কোনো ভিত্তি অজানা ছিল না।
